\documentclass[times, utf8, diplomski]{fer}
\usepackage{booktabs}

\begin{document}

\thesisnumber{1472}

\title{Multimodalno učenje manipulacijskih akcija iz vizualnih, jezičnih i taktilnih
podataka}

\author{Fran Maric}

\maketitle

% Ispis stranice s napomenom o umetanju izvornika rada. Uklonite naredbu \izvornik ako želite izbaciti tu stranicu.
\izvornik

\zahvala{Najvise se zelim zahvaliti mojim roditeljima Miri i Franji na njihovom odgoju, cijelo zivotnoj podrsci i pomoci. Zahvaljujem se mojoj djevojci Eli na razumijevanju, podrsci i motivaciji u najbitnijim trenucima. Zatim se zahvaljujem svojoj sestri Ivi koja je bila velika podrska. Zelim se zahvaliti mom prijatelju Hrvoju koji mi je pomogao pri motivaciji i ispunjavanju mojih rupa u znanju znanosti o podatcima. Velika zahvala ide direktno svim ljudima iz Larics laboratorija bez cije podrske ovaj rad ne bi bio moguc. Od njih zelim istaknuti asistenta Peru Drobac, asistenta Brunu Marica i mentora Matka Orsaga.}

\tableofcontents

\chapter{Uvod}
Podrucje robotike sve vise se oslanja na neuronske mreze. Neuronske mreze su se pokazale kao vrlo efektivne za obradu sliku, ekstrakciju semantickog znacenja iz slika, obradu jezika i slicne primjene. Transformer je nova arhitektura neuronskih mreza koja se pokazala vrlo znacajnom u obradi prirodnog jezika. Takoder je znacajno utjecala na podrucje obrade slike. Ti napretci u zadnjih pet godina su motivirali istrazivace da primjene istu arhitekturu neuronskih mreza na zadatke manipulacije robotskim rukama. U ovom radu cemo pregledati podrucje neuronskih mreza baziranih na transformer arhitekturama koje su ostavile znacajan utjecaj na podrucje robotike. Izradit cemo podatke temeljene na ljudskim demonstracijama i evaluirati neuronsku mrezu treniranu na tim podatcima.

\chapter{Motivacija}

\section{Zadatak}
moj zadatak je bio da budem jako marljivi uporan
\section{Hipoteza}

\chapter{Pregled podrucja}
\section{LLM}
\section{VLM}
\section{VLA}
\subsection{RT2}
\subsection{OpenVLA}
\subsection{Gr00t}
\subsection{Pi0}
\section{VLA sa silom}
\subsection{ForceVLA}

\chapter{Postava}
\section{Senzor sile}
sirovi podatci su vrlo sumoviti i driftaju, EMA filter
\section{Kamere}
web camera, 3D printani mount
\section{Kalup}
- 3D printani kako ne bi prezasitio senzor sile
\section{ROS}

\chapter{Demonstracija}
sinkronizacija podataka, 30Hz, vecina koristi teleoperaciju dok mi koristimo slobodne pokrete uhvacene pomocu optitracka, izvodenje pokreta

\chapter{Ucenje}
\section{Arhitektura mreze}
\section{Reprezentacija akcije}
- joint velocity
- delta EE pose
- absolute EE pose
\section{Gradijentni spust}

\chapter{Evaluacija}
\section{Metrika}
- sila:
    - konstantna sila tijekom kontakta?
    - prosjecna sila manja od thresholda
    - postotak vremena kada sila prelazi F\_max
    - glatkoca trajektorije - jerk derivacija akceleracije
    - konzistentnost brzine
- generalizcija:
    - novi kalup
    - nova orijentacija kalupa?
    - promjena jezicne upute (poliraj vs brusi vs izgladi vs obradi povrsinu)
\section{Ablacije}
- samo Z sila ili sve sile i torqueovi

\chapter{Zaključak}

\bibliography{literatura}
\bibliographystyle{fer}

\begin{sazetak}
Sažetak na hrvatskom jeziku.

\kljucnerijeci{VLA, sila, manipulacija}
\end{sazetak}

\engtitle{Multimodal learning of manipulation actions from visual, linguistic and force data}
\begin{abstract}
Abstract.

\keywords{VLA, force, manipulation}
\end{abstract}

\end{document}